\color{unidarkgreen}{\subsection{Максимальная токовая защита}}
\color{black}

\begin{enumerate}[label=\arabic{section}.\arabic{subsection}.\arabic*,  align=left, labelsep=4pt, leftmargin=0pt, itemindent=57pt]

\item
Максимальная токовая защита с комбинированным пуском по напряжению предназначена для
резервирования основных защит автотрансформатора и резервирования отключения КЗ на шинах.

\begin{figure}[H]
\centering
\includegraphics[width=0.7\textwidth,height=0.7\textheight,keepaspectratio]{\fbpath/661.0301 МТЗ/img/MTZ.pdf}
\captionof{figure}{Функциональная схема МТЗ-1}\label{pict_mtz:MTZ}
\end{figure}

\item
В составе логики максимальной токовой защиты предусмотрено две ступени МТЗ, ввод в работу которых
осуществляется с помощью программных переключателей \textbf{«МТЗ-1(2)»}. С помощью уставки \textbf{«ИО МТЗ»} имеется
возможность выбрать по фазным или линейным токам (по разности токов) будут работать измерительные
органы МТЗ.

\item
Ток срабатывания ступеней МТЗ задается с помощью уставок \textbf{«Iср МТЗ-1(2)»}, а выдержки времени
срабатывания ступеней регулируются уставками \textbf{«Тср МТЗ-1(2)»}.

\item
Имеется возможность задать различные режимы работы ступеней МТЗ при возникновении
неисправности в цепях напряжения. Вариант действия логики направленных ступеней МТЗ при срабатывании
БНН задается программными переключателями \textbf{«Режим БНН МТЗ-1(2)»}, имеющими три положения:
\begin{itemize}
\item \textit{«Выведено»} -- режим работы ступени не зависит от БНН;
\item \textit{«ОНМ»} – ступень переводится в ненаправленный режим работы;
\item \textit{«МТЗ»} – при срабатывании БНН ступень выводится из работы. Однако, если ступень задана ненаправленной, то она останется в работе независимо от срабатывания БНН. 
\end{itemize}

\item
Для обеспечения несрабатывания МТЗ от бросков намагничивающего тока предусмотрена блокировка
ступеней МТЗ по содержанию второй гармоники в фазном или линейном токе в зависимости от уставки \textbf{«ИО
МТЗ»}. Блокировка основана на контроле отношения действующего значения второй гармоники тока к основной
гармонике. Ступени МТЗ блокируется при превышении данного соотношения пороговой величины, заданной
уставкой \textbf{«I2г/I1г МТЗ»}. Наличие функции блокировки от бросков намагничивающего тока задается для
каждой ступени отдельно с помощью программных переключателей \textbf{«Режим БНТ МТЗ-1(2)»}. Для
предотвращения излишней блокировки при больших токах КЗ, сопровождающихся насыщением ТТ,
блокировка выводится из действия при превышении тока  основной гармоники выше уставки \textbf{«Iблок ОВ БНТ
МТЗ»}.

\item
Для повышения чувствительности, предусмотрена возможность пуска ступени МТЗ от
комбинированного пускового органа напряжения первой системы шин (КПОН 1СШ), которая задается
программными переключателями \textbf{«Режим КПОН 1СШ МТЗ-1(2)»} и от комбинированного пускового органа
напряжения второй системы шин (КПОН 2СШ), которая задается программными переключателями \textbf{«Режим
КПОН 2СШ МТЗ-1(2)»}.

\item
Направленность МТЗ обеспечивается органами направления мощности выполненными по 90-градусной
схеме с парами токов и напряжений Ia и Ubc, Ib и Uca, Ic и Uab. Направление мощности определяется по
величине фазового угла между током и напряжением отдельно для каждой пары сигналов. Разрешение работы
от органа направленности МТЗ будет происходить от пары сигналов с максимальным током. За положительное
направление векторов принято вращение против часовой стрелки. Углы максимальной чувствительности ОНМ
откладываются от линейных напряжений и задаются уставками \textbf{«Фмч ОНМ в 1СШ»} и \textbf{«Фмч ОНМ в 2СШ»}.
Уставки по углам максимальной чувствительности рекомендуется задавать для ОНМ в 1СШ -- 45º, а для ОНМ
в 2СШ -- 225º. Ширина зоны срабатывания ± 85º от угла максимальной чувствительности. Величина тока работы
ОНМ 0,05$\cdot I_{\textsl{ном}}$, величина напряжения работы ОНМ 0,01$\cdot U_{\textsl{ном}}$.

\begin{figure}[H]
\centering
\includegraphics[width=1\textwidth,height=1\textheight,keepaspectratio]{\fbpath/661.0301 МТЗ/img/ONM_MTZ.pdf}
\captionof{figure}{Функциональная схема ОНМ МТЗ}\label{pict_mtz:ONM_MTZ}
\end{figure}

\item
Для выбора режима направленности по мощности предусмотрены программные переключатели
\textbf{«Направл МТЗ-1(2)»}, имеющий три положения:
\begin{itemize}
\item \textit{«Ненапр»} -- ступень выполнена ненаправленной;
\item \textit{«В 1СШ»} -- ступень выполнена направленной в сторону 1СШ;
\item \textit{«В 2СШ»} -- ступень выполнена направленной в сторону 2СШ.
\end{itemize}

\item
Оперативный ввод/вывод МТЗ производится с помощью БУ (блоков управления функциями) \textit{«Вывод МТЗ»}. 

\item
В составе логики максимальной токовой защиты предусмотрена логика оперативного ускорения и логика
пуска автоматического ускорения ступеней МТЗ.

\begin{figure}[H]
\centering
\includegraphics[width=0.7\textwidth,height=0.7\textheight,keepaspectratio]{\fbpath/661.0301 МТЗ/img/GL_MTZ.pdf}
\captionof{figure}{Функциональная схема общей логики МТЗ}\label{pict_mtz:GL_MTZ}
\end{figure}

\item
Оперативное ускорение МТЗ предназначено для оперативного ускорения ступени МТЗ при выводе ДЗШ.
Ввод в работу ОУ МТЗ и выбор ускоряемой ступени задается программным переключателем \textbf{«ОУ МТЗ»}.
Предусмотрена возможность вывода направленности ускоряемой ступени с помощью программного ключа
\textbf{«Направл МТЗ при ОУ»}. Действие ускорения осуществляется с регулируемой независимой выдержкой
времени \textbf{«Тср ОУ МТЗ»}. При выводе направленности ускоряемой ступени ее блокировка от БНТ
осуществляется с помощью программного ключа \textbf{«Режим БНТ ОУ МТЗ ненапр»}.

\item
Оперативный ввод ОУ МТЗ производится с помощью команды управления \textit{«Ввод ОУ МТЗ»}.

\item
Ввод в работу пуска АУ МТЗ и выбор ускоряемой ступени задается программным переключателем \textbf{«АУ
МТЗ»}. Предусмотрена возможность вывода направленности ускоряемой ступени с помощью программного
ключа \textbf{«Направл МТЗ при АУ»}. При выводе направленности ускоряемой ступени ее блокировка от БНТ
осуществляется с помощью программного ключа \textbf{«Режим БНТ АУ МТЗ ненапр»}.

\small
\setlength{\extrarowheight}{0.05cm} % высота ячеек в таблице
\begin{longtable}
{|>{\centering\arraybackslash}m{4.5cm}|>{\centering\arraybackslash}m{4.5cm}|>{\centering\arraybackslash}m{3cm}|>{\centering\arraybackslash}m{3.5cm}|}
\caption{Перечень уставочных параметров МТЗ\hfill\vspace{-0.5\baselineskip}}\label{mtz:tbl1} \\ 
\hhline{|-|-|-|-|}
\rowcolor{gray!30}
Наименование уставки & Диапазон & Шаг &  Значение по умолчанию \\ 
\hhline{|-|-|-|-|}
\endfirsthead

\caption*{\hspace{3pt}\emph{Продолжение таблицы \ref{mtz:tbl1}\hfill\vspace{-0.5\baselineskip}}} \\ % сделано по ГОСТ 2.105 п.6.8.7
\hhline{|-|-|-|-|}
\rowcolor{gray!30}
Наименование уставки & Диапазон & Шаг &  Значение по умолчанию \\  
\hhline{|-|-|-|-|}
\endhead
\hhline{|-|-|-|-|}

\endfoot
\endlastfoot

%===m>МТЗ|-
%===m

\end{longtable}
\setlength{\extrarowheight}{0.0cm} % возврат параметра высота ячеек в таблице
\normalsize





























































\end{enumerate} 